\subsection{INDICATRIX OF A COMPOSITION OPERATOR}

With the same hypotheses as in \(n^{\circ}\) 5, the formula (26) of (VI, p.~277) applied to the function \(e^{\lambda x}\) , gives

\[
U _ {x} ^ {\xi} (e ^ {\lambda \xi}) = U _ {0} ^ {\xi} (e ^ {\lambda x} e ^ {\lambda \xi}) = e ^ {\lambda x} U _ {0} ^ {\xi} (e ^ {\lambda \xi}) = u (\lambda) e ^ {\lambda x}\tag{30}
\]

on putting \(u(\lambda) = U_{0}^{\xi}(e^{\lambda\xi})\) . One says that the function \(u(\lambda)\) , defined on C, with complex values, is the indicatrix of the composition operator U. Note that if the restriction of U to the ring C{[}X{]} of polynomials is equal to the series

\[
\mathbf {D} ^ {p} \sum_ {n = 0} ^ {\infty} \alpha_ {n} \mathbf {D} ^ {n}
\]

\pandocbounded{\includegraphics[keepaspectratio]{images/part002_816cd56c.jpg}}

(VI, p.~271, th. 1) (which we have denoted \(u(\mathrm{D})\) in VI, p.~271), the series of complex terms with general term \(\alpha_{n}\lambda^{n+p}\) is not necessarily convergent for \(\lambda \neq 0\) , and that even if it converges for certain values of \(\lambda\) , the sum is not necessarily equal to the indicatrix \(u(\lambda)\) of U (VI, p.~291, exerc. 2). We shall say that the composition operator U is regular if there exists a neighbourhood of 0 in C such that the series with general term \(\alpha_{n}\lambda^{n+p}\) is absolutely convergent and has sum equal to the indicatrix \(u(\lambda)\) on this neighbourhood \(^{1}\) . Let us apply formula (27) of VI, p.~278, to the function \(e^{\lambda x}\) , making h = 0; since \(\mathrm{D}^{m}(e^{\lambda x}) = \lambda^{m} e^{\lambda x}\) one has \(U_{0}^{\xi}\left(\mathrm{D}^{m}(e^{\lambda x})\right) = \lambda^{m} u(\lambda)\) ; it then follows that for every complex \(\lambda\) such that \(u(\lambda) \neq 0\)

\[
\frac {\lambda^ {p} e ^ {\lambda x}}{u (\lambda)} = \sum_ {m = 0} ^ {n} u _ {m} (x) \frac {\lambda^ {m}}{m !} - \frac {\lambda^ {n + 1}}{u (\lambda)} U _ {0} ^ {\xi} \left(\int_ {0} ^ {\xi - x} \frac {1}{n !} u _ {n} (x + \eta) e ^ {\lambda (\xi - \eta)} d \eta\right)\tag{31}
\]

and in particular, for \(x = 0\)

\[
\frac {\lambda^ {p}}{u (\lambda)} = \sum_ {m = 0} ^ {n} \beta_ {m} \frac {\lambda^ {m}}{m !} - \frac {\lambda^ {n + 1}}{u (\lambda)} U _ {0} ^ {\xi} \left(\int_ {0} ^ {\xi} \frac {1}{n !} u _ {n} (\eta) e ^ {\lambda (\xi - \eta)} d \eta\right)\tag{32}
\]

with \(\beta_{m} = u_{m}(0)\) .

If \(U\) is a regular operator, for all \(\lambda \in \mathbf{C}\) such that the entire series \(u(\lambda) = \sum_{n=0}^{\infty} \alpha_n \lambda^{n+p}\) and \(\sum_{n=0}^{\infty} \beta_n \frac{\lambda^n}{n!}\) are absolutely convergent \(^2\) , it follows from formula (16) and the formula for the product of two absolutely convergent series (Gen.~Top., VIII, p.~115, prop. 1) that one has

\[
{\frac {\lambda^ {p}}{u (\lambda)}} = \sum_ {n = 0} ^ {\infty} \beta_ {n} {\frac {\lambda^ {n}}{n !}}.\tag{33}
\]

Similarly, since the Taylor expansion of \(e^{\lambda x}\) is absolutely convergent for every \(\lambda \in \mathbf{C}\) and every \(x \in \mathbf{C}\) (III, p.~106) one also has (formulae (6) (VI, p.~273) and (16) (VI, p.~274)), for all the values considered and for all \(x \in \mathbf{C}\)

\[
\frac {\lambda^ {p} e ^ {\lambda x}}{u (\lambda)} = \sum_ {n = 0} ^ {\infty} u _ {n} (x) \frac {\lambda^ {n}}{n !}.\tag{34}
\]

Remark. One can use formula (33) (resp. (34)) to calculate the \(\beta_{n}\) (resp. the \(u_{n}(x)\) ) by using the following lemma on entire series:

Lemma. If two entire series \(\sum_{n=0}^{\infty} c_n \lambda^n\) , \(\sum_{n=0}^{\infty} d_n \lambda^n\) are absolutely convergent for all \(\lambda\) in a neighbourhood of 0, and if \(\sum_{n=0}^{\infty} c_n \lambda^n = \sum_{n=0}^{\infty} d_n \lambda^n\) for these values of \(\lambda\) , then \(c_n = d_n\) for all integers \(n \geqslant 0^3\) .

If, by any procedure, one can find an entire series that converges to \(\lambda^p / u(\lambda)\) on a neighbourhood of 0, then the coefficients of this series are necessarily equal to the \(\beta_n / n!\) . It is this procedure that we shall apply in the examples that follow.

Examples. 1) If \(U\) is the identity map one has \(u(\lambda) = 1\) and the operator \(U\) is clearly regular; since \(u_{n}(x) = x^{n}\) the formula (27) of VI, p.~278, can be written, putting \(t = \xi - \eta\)

\[
f (x + h) = \sum_ {n = 0} ^ {\infty} \frac {1}{m !} f ^ {(m)} (h) x ^ {m} + \int_ {h} ^ {x + h} f ^ {(n + 1)} (t) \frac {(x + h - t) ^ {n}}{n !} d t
\]

that is, it reduces to the Taylor formula (II, p.~62).

\begin{enumerate}
\def\labelenumi{\arabic{enumi})}
\setcounter{enumi}{1}
\tightlist
\item
  Let us take for U the composition operator which, to every function f defined on \(R_{+}\) , associates the function \(x \mapsto f(x + 1) - f(x)\) ; then
\end{enumerate}

\[
U _ {x} ^ {\xi} (f (\xi)) = f (x + 1) - f (x);
\]

we have seen (VI, p.~275) that the restriction of U to C{[}X{]} is equal to \(e^{D}-1\) . Since, on the other hand, \(u(\lambda)=e^{\lambda}-1\) , the operator U is regular; we shall see (VI, p.~288) how to determine the Bernoulli numbers \(b_{n}\) by calculating an expansion of \(\frac{\lambda}{e^{\lambda}-1}\) as a convergent entire series. On applying formula (27) of VI, p.~278 to a primitive of a function f admitting a continuous \(n^{th}\) derivative on \(R_{+}\) , we obtain

\[
\begin{array}{l} f (x + h) = \int_ {h} ^ {h + 1} f (t) d t \\ \qquad + \sum_ {m = 1} ^ {n} \frac {1}{m !} \mathrm{B} _ {m} (x) \left(f ^ {(m - 1)} (h + 1) - f ^ {(m - 1)} (h)\right) + \mathrm{R} _ {n} (x, h) \end{array}\tag{35}
\]

\[
\begin{array}{l} \mathrm{R} _ {n} (x, h) = - \int_ {0} ^ {1 - x} \frac {\mathrm{B} _ {n} (x + \eta)}{n !} f ^ {(n)} (h + 1 - \eta) d \eta \\ \qquad \qquad + \int_ {0} ^ {- x} \frac {\mathrm{B} _ {n} (x + \eta)}{n !} f ^ {(n)} (h - \eta) d \eta . \end{array}\tag{36}
\]

*3) Let \(\mathbf{E}\) be the vector space of functions \(f\) defined and continuous on \(\mathbf{R}\) , and such that the integral \(\int_{-\infty}^{+\infty} f(x + \xi)e^{-\xi^2 / 2}d\xi\) converges for all \(x \geqslant 0\) . The operator \(\mathrm{U}\) defined by

\[
U _ {x} ^ {\xi} (f (\xi)) = \frac {1}{\sqrt {2 \pi}} \int_ {- \infty} ^ {+ \infty} e ^ {- \xi^ {2} / 2} f (x + \xi) d \xi
\]

is then defined on E and is clearly a composition operator. The space E contains all the exponentials \(e^{\lambda x}\) ( \(\lambda\) arbitrary complex), and one has

\[
\begin{array}{l} u (\lambda) = \frac {1}{\sqrt {2 \pi}} \int_ {- \infty} ^ {+ \infty} e ^ {- (\xi^ {2} / 2) + \lambda \xi} d \xi \\ = \frac {1}{\sqrt {2 \pi}} e ^ {\lambda^ {2} / 2} \int_ {- \infty} ^ {+ \infty} e ^ {- (\xi - \lambda) ^ {2} / 2} d \xi = e ^ {\lambda^ {2} / 2} \end{array}
\]

(cf.~III, p.~120, exerc. 24, and VII, p.~313, formula (22)). One has \(n! \alpha_n = U_0^\xi (\xi^n) = \frac{1}{\sqrt{2\pi}} \int_{-\infty}^{+\infty} e^{-\xi^2 / 2} \xi^n d\xi\) . For every integer \(n\) one can write

\[
\sum_ {k = 0} ^ {n} \int_ {- \infty} ^ {+ \infty} \frac {| \lambda \xi | ^ {k}}{k !} e ^ {- \xi^ {2} / 2} d \xi \leqslant 2 \int_ {0} ^ {+ \infty} e ^ {- (\xi^ {2} / 2) + | \lambda | \xi} d \xi .
\]

The series \(\sum_{n=0}^{\infty}e^{-\xi^{2}/2}\frac{(\lambda\xi)^{n}}{n!}\) can therefore be integrated term-by-term over R (II, p.~72, cor. 1), which proves that the series \(\sum_{n=0}^{\infty}\alpha_{n}\lambda^{n}\) converges absolutely for every \(\lambda\in C\) , and has a sum equal to \(u(\lambda)=e^{\lambda^{2}/2}=\sum_{n=0}^{\infty}\frac{\lambda^{2n}}{2^{n}n!}\) ; thus the operator U is regular. Applying the lemma mentioned above shows that \(\alpha_{2n}=1/2^{n}n!\) , \(\alpha_{2n+1}=0\) for every \(n\geqslant0\) ; the operator U is thus of order 0. One has

\[
\frac {1}{u (\lambda)} = e ^ {- \lambda^ {2} / 2} = \sum_ {n = 0} ^ {\infty} \frac {(- 1) ^ {n} \lambda^ {2 n}}{2 ^ {n} n !},
\]

the series being absolutely convergent for every \(\lambda \in \mathbf{C}\) ; another application of the lemma shows that \(\beta_{2n} = \frac{(-1)^n}{2^n}\frac{(2n)!}{n!}\) , \(\beta_{2n+1} = 0\) ; further, the series \(\sum_{n=0}^{\infty}\frac{\lambda^n}{n!} u_n(x)\) is absolutely convergent for every \(\lambda \in \mathbf{C}\) and every \(x \in \mathbf{R}\) , and one has

\[
\sum_ {n = 0} ^ {\infty} \frac {\lambda^ {n}}{n !} u _ {n} (x) = \exp \left(- \frac {\lambda^ {2}}{2} + \lambda x\right) = \exp \left(\frac {x ^ {2}}{2}\right) \exp \left(- \frac {1}{2} (\lambda - x) ^ {2}\right).
\]

On applying the Taylor formula to the function \(\exp(-x^2/2)\) one then obtains the following expression for the polynomials \(u_n(x)\) :

\[
u _ {n} (x) = (- 1) ^ {n} e ^ {x ^ {2} / 2} \frac {d ^ {n}}{d x ^ {n}} (e ^ {- x ^ {2} / 2}).
\]

This polynomial is called the Hermite polynomial of degree n, and is often denoted by \(\mathrm{H}_{n}(x)\) . The formulae (7), (8) and (9) of VI, p.~273, here give

\[
\frac {d \mathrm{H} _ {n}}{d x} = n \mathrm{H} _ {n - 1} (x)
\]

\[
\mathrm{H} _ {n} (x + y) = \sum_ {k = 0} ^ {n} {\binom {n} {k}} \mathrm{H} _ {n - k} (x)   y ^ {k}
\]

\[
\frac {1}{\sqrt {2 \pi}} \int_ {- \infty} ^ {+ \infty} e ^ {- \xi^ {2} / 2} \mathrm{H} _ {n} (x + \xi) d \xi = x ^ {n}
\]

and the formula (27) of VI, p.~278, becomes, for \(h = 0\)

\[
\begin{array}{l} \sqrt {2 \pi} f (x) = \sum_ {m = 0} ^ {n} \left(\int_ {- \infty} ^ {+ \infty} e ^ {- \xi^ {2} / 2} f ^ {(m)} (\xi) d \xi\right) \frac {\mathrm{H} _ {m} (x)}{m !} \\ - \int_ {- \infty} ^ {+ \infty} d \xi \int_ {0} ^ {\xi} \frac {\mathrm{H} _ {n} (x + \eta)}{n !} e ^ {- (\xi + x) ^ {2} / 2} f ^ {(n + 1)} (x + \xi - \eta) d \eta . * \end{array}
\]

